El an\'alisis global del repositorio mostr\'o que la incompletitud observacional no se distribuye de manera uniforme en el cat\'alogo. Mapper identific\'o regiones con estructura local coherente; la sombra observacional delimit\'o fronteras de detectabilidad; TOI prioriz\'o nodos; ATI prioriz\'o planetas ancla. Ese encadenamiento permiti\'o formular una hip\'otesis general de submuestreo topol\'ogico en el espacio f\'isico-orbital.

El nuevo paso metodol\'ogico consiste en preguntar si esa sospecha global puede traducirse en intervalos orbitales concretos dentro de sistemas planetarios conocidos. En vez de trabajar solo con regiones del cat\'alogo completo, el subproyecto estudia conjuntos de planetas que orbitan la misma estrella y pregunta si entre dos planetas observados existe un hueco orbital an\'omalo comparado con sistemas an\'alogos.

La idea central es simple. Si un sistema tiene planetas ordenados por per\'iodo,
\[
P_1 < P_2 < \cdots < P_n,
\]
entonces cada par adyacente define un gap. Un gap amplio no demuestra la existencia de un planeta no detectado, pero s\'i puede marcar un intervalo orbital prioritario. Si ese gap adem\'as se encuentra cerca de nodos o planetas con soporte previo de TOI, ATI o sombra observacional, el intervalo adquiere m\'as valor como candidato para seguimiento.

Este cambio de escala es importante porque conecta dos niveles del proyecto:
\begin{itemize}
\item el nivel global, donde Mapper se\~nala regiones de posible incompletitud observacional;
\item el nivel intra-sistema, donde se priorizan intervalos orbitales espec\'ificos alrededor de estrellas conocidas.
\end{itemize}

La lectura correcta del m\'odulo es de priorizaci\'on observacional. No se detectan planetas reales ni se infiere un conteo verdadero de objetos ausentes. Lo que se obtiene es un ranking prudente de candidatos a planetas no detectados e intervalos orbitales prioritarios donde nuevas observaciones podr\'ian ser m\'as informativas.
